\documentclass{article}

\usepackage{mathtools}
\usepackage{subcaption}
\usepackage{authblk}
% \usepackage{hyperref}
\usepackage{multirow}

\usepackage{wrapfig}
\usepackage{booktabs}
\usepackage{xcolor}
\usepackage{enumitem}
\usepackage{listings}
\usepackage[colorlinks=true, linkcolor=blue, citecolor=blue, urlcolor=blue]{hyperref}

\title{Laplace Analysis and Applications}

\author[1]{Pierre-Simon Laplace}
\author[2]{Oliver Heaviside}
\author[3]{Gustav Robert Kirchhoff}

\affil[1]{Acad´emie des Sciences, Paris, France}
\affil[2]{Royal Society, London, United Kingdom}
\affil[3]{University of Heidelberg, Heidelberg, Germany}

\begin{document}
\maketitle

\section{From Functions to Transforms}
\subsection{The Laplace Transform}
Let $f(t)$ be a function defined for $t \geq 0$. Its Laplace transform is defined by
\begin{equation}
    \begin{cases}
    \mathcal{L}\{f(t)\} = F(s) = \int_0^{\infty} e^{-st}f(t)dt & s>0\\
    \end{cases}
\end{equation}

The transform changes a function of time into a function of the new variable s. Symbolically,
\[
f(t) \xrightarrow{\text{L}} F(s)
\]
The inverse Laplace transform reverses this process:
\begin{equation}
    f(t) = \mathcal{L}^{-1} \{F(s)\}
\end{equation}

The main purpose of the method is to replace differentiation and integration in the time domain
by simpler algebraic operations in the transform domain. A concise overview is available from
the \href{https://mathworld.wolfram.com/LaplaceTransform.html}{Laplace Transform entry on MathWorld.}

\subsection{A Simple Example}
For the constant function $f (t) = 1$,
\begin{equation}
    \begin{aligned}
        \mathcal{L}\{1\} &= \int_0^{\infty} e^{-st} dt\\
        &= \left[
            -\frac{1}{s} e^{-st}
        \right]_0^{\infty}\\
        &= -\frac{1}{s}\\
    \end{aligned}
\end{equation}

Similarly, 
\begin{equation}
    \mathcal{L}\{t\} = \frac{1}{s^3}, \mathcal{L}\{t^2\} = \frac{2}{s^3}, 
\end{equation}
{\small The variable s is commonly interpreted as a complex-frequency variable}

\section{Common Transform Equations}
\begin{table}
    \centering
    \caption{Selected Laplace-transform equations.}
    \label{tab:transforms}
    \begin{tabular}{|c|c|c|c|}
        \hline
        \multirow{2}{*}{Family} & \multicolumn{2}{|c|}{Transform Equation} & \multirow{2}{*}{Condition} \\
        \cline{2-3}
        & Time function & Laplace transform & \\
        \hline
    \multirow{3}{*}{Algebric} &1 & $\frac{1}{s}$ & $s>0$\\
    \cline{2-4}
    &$t$ &$\frac{1}{s^2}$ & $s>0$\\
    \cline{2-4}
    &$t^n$ &$\frac{n!}{s^{n+1}}$ & $s>0, n=0,1,2,\dots$\\
    \hline
    \multirow{2}{*}{Exponential} & $e^{at}$ & $\frac{1}{s-a}$ & $s>a$ \\
    \cline{2-4}
    & $e^{-at}$ & $\frac{1}{s+a}$ & $s>-a$ \\
    \hline
    \multirow{2}{*}{Oscillatory} & $\sin(\omega t)$ & $\frac{\omega}{s^2+\omega^2}$ & $\omega>0$ \\
    \cline{2-4}
    & $\cos(\omega t)$ & $\frac{s}{s^2+\omega^2}$ & $\omega>0$ \\
    \hline
    & Unit-step function $u(t-a)$ & $\frac{e^{-as}}{s}$ & $a>0$ \\
    \hline
    \end{tabular}
\end{table}


The formulas in Table \ref{tab:transforms} allow many functions to be transformed without repeatedly evaluating improper integrals.

\section{Application to an RC Circuit}

\begin{wrapfigure}{r}{0.4\textwidth}
    \centering
    % The instructions requested a rotated figure. The 'angle=90' parameter fulfills this requirement.
    % Make sure a file named 'circuit.png' exists in your working directory.
    \includegraphics[width=\linewidth]{rc_circuit.png}
    \caption{A simple series RC circuit.}
    \label{fig:circuit}
\end{wrapfigure}

\textbf{Circuit Model} \\
Figure \ref{fig:circuit} shows a simple resistor--capacitor circuit connected to an input voltage source. The resistor limits the current, while the capacitor stores electrical energy. 

For a resistor $R$ and capacitor $C$, Kirchhoff's voltage law gives
\begin{equation}
    RC\frac{dv_{C}(t)}{dt} + v_{C}(t) = v_{in}(t), \label{eq:kvl}
\end{equation}
where $v_{C}(t)$ is the voltage across the capacitor. Assuming $v_{C}(0)=0$, taking the Laplace transform of Equation (\ref{eq:kvl}) gives
\begin{align}
    RC[sV_{C}(s) - v_{C}(0)] + V_{C}(s) &= V_{in}(s) \nonumber \\
    \Rightarrow RCsV_{C}(s) + V_{C}(s) &= V_{in}(s) \nonumber \\
    \Rightarrow V_{C}(s)(RCs + 1) &= V_{in}(s). \label{eq:laplace_rc}
\end{align}

The transfer function is therefore
\begin{equation}
    H(s) = \frac{V_{C}(s)}{V_{in}(s)} = \frac{1}{RCs+1}.
\end{equation}
For a constant input $V_0$, $V_{in}(s) = \frac{V_0}{s}$. Thus,
\begin{equation}
    V_{C}(s) = \frac{V_0}{s(RCs+1)}.
\end{equation}
Taking the inverse transform gives
\begin{equation}
    v_{C}(t) = V_0(1 - e^{-t/(RC)}). \label{eq:time_domain}
\end{equation}

The quantity $\tau = RC$ is called the time constant of the circuit. A larger value of $\tau$ causes the capacitor to charge more slowly \cite{ogata}.




\end{document}